\section*{\centering{RESUMO}}

Este trabalho investiga o desempenho de três algoritmos de \textit{Gradient Boosting} baseados em árvore --- XGBoost, LightGBM e CatBoost --- em ambiente de Aprendizado Federado aplicado à detecção de fraude em abertura de contas bancárias. Utilizando o \textit{dataset} BAF (\textit{Bank Account Fraud}) e o \textit{framework} Flower para simulação federada, foram avaliadas duas estratégias de agregação --- Bagging e Cíclica --- com 3 clientes, 5 rodadas de comunicação e particionamento IID balanceado. Os modelos federados alcançaram TPR entre 54,83\% e 56,79\% no ponto de operação de 5\% de FPR, resultados competitivos com o \textit{baseline} centralizado do BAF (40--55\%). O XGBoost Bagging obteve o melhor desempenho absoluto (56,79\% TPR, AUC 0,8934), enquanto o CatBoost Cíclico alcançou resultado praticamente equivalente (56,72\% TPR) com menor custo de comunicação (1,27 MB). A aplicação de limiares independentes por grupo demográfico elevou o FPR Ratio de $\approx$0,3 para $\approx$0,88--1,0, demonstrando que a equidade algorítmica pode ser alcançada com perda média de apenas 2,0 pontos percentuais em TPR. Os resultados demonstram a viabilidade de modelos baseados em árvore para sistemas federados de detecção de fraude, oferecendo eficiência de comunicação superior a abordagens baseadas em redes neurais profundas.

{\bf Palavras-chave:} Aprendizado Federado, Detecção de Fraudes Bancárias, XGBoost, LightGBM, CatBoost, Flower Framework, Privacidade de Dados, Modelos Baseados em Árvore.
